\begin{bt}%[9D6H18-2]
% nguon: Tuyển tập đề tham khảo môn Toán Tuyển sinh 10 năm học 2025-2026 TP.HCM – Đề tham khảo Quận Bình Thạnh (Đề 1) – Bài 1
Cho hàm số $y=\dfrac{1}{4}x^2$.
\begin{enumerate}[a)]
	\item Vẽ đồ thị $(P)$ của hàm số trên.
	\item Tìm các điểm $A, B$ thuộc $(P)$ sao cho tam giác $OAB$ vuông cân tại $O$.
\end{enumerate}
\loigiai{
\begin{enumerate}[a)]
	\item Bảng giá trị tương ứng của $x$ và $y$:
	\begin{center}
	\begin{tabular}{|c|c|c|c|c|c|}
	\hline
	$x$ & $-4$ & $-2$ & $0$ & $2$ & $4$ \\ \hline
	$y=\dfrac{1}{4}x^2$ & $4$ & $1$ & $0$ & $1$ & $4$ \\ \hline
	\end{tabular}
	\end{center}
	Đồ thị $(P)$ là một parabol có đỉnh là gốc toạ độ $O(0;0)$, nhận trục tung $Oy$ làm trục đối xứng và quay bề lõm lên phía trên.
	\item Giả sử hai điểm $A(x_1; y_1)$ và $B(x_2; y_2)$ thuộc $(P)$ khác $O$. Khi đó $y_1=\dfrac{1}{4}x_1^2 > 0$ và $y_2=\dfrac{1}{4}x_2^2 > 0$.\\
	Do $(P)$ đối xứng qua trục tung $Oy$, tam giác $OAB$ vuông cân tại $O$ khi và chỉ khi hai điểm $A$ và $B$ đối xứng nhau qua trục tung $Oy$ và đường thẳng $OA, OB$ tạo với trục hoành một góc $45^\circ$, tức là:
	$$
	x_2 = -x_1 \quad \text{và} \quad y_1 = y_2 = |x_1|.
	$$
	Thay vào phương trình parabol $(P)$, ta có:
	$$
	\dfrac{1}{4}x_1^2 = |x_1| \iff |x_1| \left( \dfrac{1}{4}|x_1| - 1 \right) = 0.
	$$
	Vì $A \ne O$ nên $|x_1| = 4 \iff x_1 = \pm 4$.\\
	Với $x_1 = 4 \implies y_1 = 4$, khi đó $x_2 = -4 \implies y_2 = 4$.\\
	Vậy hai điểm cần tìm là $A(4;4)$ và $B(-4;4)$ (hoặc hoán vị $A(-4;4)$ và $B(4;4)$).
\end{enumerate}
}
\end{bt}

\begin{bt}%[9D6H20-2]
% nguon: Tuyển tập đề tham khảo môn Toán Tuyển sinh 10 năm học 2025-2026 TP.HCM – Đề tham khảo Quận Bình Thạnh (Đề 1) – Bài 2
Cho phương trình bậc hai: $2x^2-3x-4=0$.
\begin{enumerate}[a)]
	\item Chứng minh phương trình trên luôn có hai nghiệm phân biệt.
	\item Không giải phương trình, hãy tính giá trị của biểu thức:
	$$
	T = \dfrac{x_1+2}{x_2} + \dfrac{x_2+2}{x_1}.
	$$
\end{enumerate}
\loigiai{
\begin{enumerate}[a)]
	\item Phương trình $2x^2-3x-4=0$ có các hệ số $a=2$, $b=-3$, $c=-4$.\\
	Vì $a \cdot c = 2 \cdot (-4) = -8 < 0$ nên phương trình luôn có hai nghiệm phân biệt $x_1, x_2$.
	\item Theo hệ thức Vi-ét, ta có:
	$$
	\begin{cases}
	S = x_1+x_2 = \dfrac{3}{2} \\
	P = x_1 x_2 = \dfrac{-4}{2} = -2.
	\end{cases}
	$$
	Biến đổi biểu thức $T$:
	$$
	T = \dfrac{x_1(x_1+2) + x_2(x_2+2)}{x_1 x_2} = \dfrac{x_1^2 + x_2^2 + 2(x_1+x_2)}{x_1 x_2} = \dfrac{(x_1+x_2)^2 - 2x_1 x_2 + 2(x_1+x_2)}{x_1 x_2}.
	$$
	Thay $S = \dfrac{3}{2}$ và $P = -2$ vào biểu thức $T$, ta được:
	$$
	T = \dfrac{\left(\dfrac{3}{2}\right)^2 - 2(-2) + 2 \cdot \dfrac{3}{2}}{-2} = \dfrac{\dfrac{9}{4} + 4 + 3}{-2} = \dfrac{\dfrac{37}{4}}{-2} = -\dfrac{37}{8} = -4{,}625.
	$$
\end{enumerate}
}
\end{bt}

\begin{bt}%[9D8H26-3]
% nguon: Tuyển tập đề tham khảo môn Toán Tuyển sinh 10 năm học 2025-2026 TP.HCM – Đề tham khảo Quận Bình Thạnh (Đề 1) – Bài 3
Gieo $3$ con súc sắc cân đối, kết quả mỗi lần gieo là một bộ thứ tự $(x, y, z)$ với $x, y, z$ lần lượt là số chấm xuất hiện trên mỗi con súc sắc. Thực hiện gieo $20$ lần, người ta thu được bảng số liệu sau:
\begin{center}
\begin{tabular}{|c|c|c|c|c|}
\hline
$(1; 4; 2)$ & $(4; 6; 2)$ & $(6; 5; 3)$ & $(2; 1; 2)$ & $(4; 4; 6)$ \\ \hline
$(4; 2; 3)$ & $(2; 4; 4)$ & $(1; 2; 1)$ & $(5; 5; 6)$ & $(5; 1; 5)$ \\ \hline
$(4; 6; 6)$ & $(3; 3; 4)$ & $(1; 1; 5)$ & $(2; 5; 4)$ & $(4; 1; 6)$ \\ \hline
$(1; 1; 4)$ & $(1; 6; 4)$ & $(2; 2; 2)$ & $(6; 2; 3)$ & $(5; 4; 6)$ \\ \hline
\end{tabular}
\end{center}
\begin{enumerate}[a)]
	\item Trong $20$ lần gieo trên, tổng số chấm lớn nhất và nhỏ nhất lần lượt là bao nhiêu? Có bao nhiêu trường hợp tích số chấm chia hết cho $2$ và đồng thời tổng số chấm là một số nguyên tố?
	\item Tính xác suất thực nghiệm của biến cố: ``Tổng số chấm lớn hơn hoặc bằng $3$ và không vượt quá $11$''.
\end{enumerate}
\loigiai{
\begin{enumerate}[a)]
	\item Tính tổng số chấm $S = x+y+z$ cho từng bộ kết quả:
	\begin{itemize}
		\item Hàng 1: $7; 12; 14; 5; 14$.
		\item Hàng 2: $9; 10; 4; 16; 11$.
		\item Hàng 3: $16; 10; 7; 11; 11$.
		\item Hàng 4: $6; 11; 6; 11; 15$.
	\end{itemize}
	Tổng số chấm lớn nhất là $16$ (xuất hiện ở các bộ $(5; 5; 6)$ và $(4; 6; 6)$).\\
	Tổng số chấm nhỏ nhất là $4$ (xuất hiện ở bộ $(1; 2; 1)$).\\
	Tích số chấm $x \cdot y \cdot z$ chia hết cho $2$ khi và chỉ khi trong bộ có ít nhất một số chẵn. Các số nguyên tố có thể xuất hiện từ tổng là $5, 7, 11$.\\
	Các kết quả thoả mãn cả hai điều kiện (có ít nhất một số chẵn và tổng là số nguyên tố) là:
	$$
	(1; 4; 2) \ [S=7], \quad (2; 1; 2) \ [S=5], \quad (2; 5; 4) \ [S=11], \quad (4; 1; 6) \ [S=11], \quad (1; 6; 4) \ [S=11], \quad (6; 2; 3) \ [S=11].
	$$
	Như vậy có tất cả $6$ trường hợp thoả mãn.
	\item Các kết quả có tổng số chấm thoả mãn $3 \le S \le 11$ gồm có $14$ trường hợp:
	$$
	(1; 4; 2), (2; 1; 2), (4; 2; 3), (2; 4; 4), (1; 2; 1), (5; 1; 5), (3; 3; 4), (1; 1; 5), (2; 5; 4), (4; 1; 6), (1; 1; 4), (1; 6; 4), (2; 2; 2), (6; 2; 3).
	$$
	Xác suất thực nghiệm của biến cố cần tìm là:
	$$
	P = \dfrac{14}{20} = \dfrac{7}{10} = 0{,}7 \ (70\%).
	$$
\end{enumerate}
}
\end{bt}

\begin{bt}%[9D6V21-3]
% nguon: Tuyển tập đề tham khảo môn Toán Tuyển sinh 10 năm học 2025-2026 TP.HCM – Đề tham khảo Quận Bình Thạnh (Đề 1) – Bài 4
Một khu vườn hình chữ nhật có chiều dài $30\text{ m}$ và chiều rộng $10\text{ m}$. Người ta làm hai lối đi trong khu vườn có dạng hình chữ nhật vuông góc với nhau và cùng có bề rộng là $x\text{ (m)}$ ($0 < x < 10$).
\begin{enumerate}[a)]
	\item Viết biểu thức $S$ biểu diễn theo $x$ diện tích của khu vườn còn lại sau khi mở lối đi.
	\item Biết rằng sau khi làm lối đi thì diện tích của khu vườn giảm đi $76\text{ m}^2$. Tìm giá trị của $x$ (làm tròn đến hàng phần mười của mét).
\end{enumerate}
\loigiai{
\begin{enumerate}[a)]
	\item Ban đầu diện tích của khu vườn là:
	$$
	S_0 = 30 \cdot 10 = 300\text{ (m}^2\text{)}.
	$$
	Khi làm hai lối đi vuông góc cắt nhau cùng có bề rộng $x\text{ (m)}$, phần diện tích đất còn lại có thể ghép lại thành một hình chữ nhật có kích thước là $(30-x)\text{ m}$ và $(10-x)\text{ m}$.\\
	Do đó biểu thức diện tích khu vườn còn lại là:
	$$
	S(x) = (30-x)(10-x) = x^2 - 40x + 300\text{ (m}^2\text{)}.
	$$
	\item Diện tích khu vườn giảm đi chính là tổng diện tích của hai lối đi trừ phần giao nhau:
	$$
	S_{\text{lối đi}} = 30x + 10x - x^2 = 40x - x^2\text{ (m}^2\text{)}.
	$$
	Theo đề bài, diện tích khu vườn giảm đi $76\text{ m}^2$, ta có phương trình:
	$$
	40x - x^2 = 76 \iff x^2 - 40x + 76 = 0.
	$$
	Giải phương trình bậc hai:
	$$
	\Delta' = (-20)^2 - 1 \cdot 76 = 400 - 76 = 324 = 18^2.
	$$
	Phương trình có hai nghiệm:
	$$
	x_1 = 20 - 18 = 2\text{ (nhận do } 0 < x < 10\text{)}, \quad x_2 = 20 + 18 = 38\text{ (loại)}.
	$$
	Vậy bề rộng của lối đi là $x = 2\text{ m}$.
\end{enumerate}
}
\end{bt}

\begin{bt}%[9H0V31-3]
% nguon: Tuyển tập đề tham khảo môn Toán Tuyển sinh 10 năm học 2025-2026 TP.HCM – Đề tham khảo Quận Bình Thạnh (Đề 1) – Bài 5
Một thúng gạo được vun đầy có phần dưới dạng nửa hình cầu với đường kính $50\text{ cm}$, phần gạo vun lên phía trên có dạng hình nón với chiều cao $15\text{ cm}$ và bán kính đáy bằng bán kính miệng thúng.
\begin{enumerate}[a)]
	\item Tính thể tích phần gạo chứa trong thúng (kết quả làm tròn đến hàng phần mười của $\text{cm}^3$).
	\item Gia đình bạn An dùng một lon sữa bò hình trụ có bán kính đáy $4\text{ cm}$, chiều cao $10\text{ cm}$ để đong gạo mỗi ngày. Biết mỗi ngày gia đình ăn hết $5$ lon gạo và mỗi lần đong thì lượng gạo chiếm $90\%$ thể tích lon. Hỏi với lượng gạo trong thúng trên thì gia đình bạn An có thể ăn nhiều nhất trong bao nhiêu ngày?
\end{enumerate}
\loigiai{
\begin{enumerate}[a)]
	\item Bán kính của thúng gạo là $R = \dfrac{50}{2} = 25\text{ cm}$.\\
	Thể tích phần nón của gạo vun lên là:
	$$
	V_{\text{nón}} = \dfrac{1}{3}\pi R^2 h = \dfrac{1}{3}\pi \cdot 25^2 \cdot 15 = 3125\pi\text{ (cm}^3\text{)}.
	$$
	Thể tích phần nửa hình cầu phía dưới của thúng là:
	$$
	V_{\text{bán cầu}} = \dfrac{2}{3}\pi R^3 = \dfrac{2}{3}\pi \cdot 25^3 = \dfrac{31250}{3}\pi\text{ (cm}^3\text{)}.
	$$
	Tổng thể tích phần gạo trong thúng là:
	$$
	V_{\text{gạo}} = V_{\text{nón}} + V_{\text{bán cầu}} = 3125\pi + \dfrac{31250}{3}\pi = \dfrac{40625}{3}\pi \approx 42542{,}4\text{ cm}^3.
	$$
	\item Thể tích của một lon sữa bò hình trụ là:
	$$
	V_{\text{lon}} = \pi r^2 h_t = \pi \cdot 4^2 \cdot 10 = 160\pi\text{ (cm}^3\text{)}.
	$$
	Mỗi lần đong gạo chiếm $90\%$ thể tích lon, do đó thể tích gạo trong mỗi lon là:
	$$
	V_{1} = 0{,}9 \cdot 160\pi = 144\pi\text{ (cm}^3\text{)}.
	$$
	Mỗi ngày gia đình ăn $5$ lon gạo, thể tích gạo tiêu thụ mỗi ngày là:
	$$
	V_{\text{ngày}} = 5 \cdot 144\pi = 720\pi\text{ (cm}^3\text{)}.
	$$
	Số ngày gia đình bạn An có thể ăn được lượng gạo trên là:
	$$
	n = \dfrac{V_{\text{gạo}}}{V_{\text{ngày}}} = \dfrac{\dfrac{40625}{3}\pi}{720\pi} = \dfrac{40625}{2160} = \dfrac{8125}{432} \approx 18{,}8\text{ (ngày)}.
	$$
	Vì lượng gạo chỉ đủ cho $18$ ngày ăn trọn vẹn (ngày thứ $19$ không đủ $5$ lon), nên gia đình bạn An có thể ăn nhiều nhất là $18$ ngày.
\end{enumerate}
}
\end{bt}

\begin{bt}%[9D1H3-4]
% nguon: Tuyển tập đề tham khảo môn Toán Tuyển sinh 10 năm học 2025-2026 TP.HCM – Đề tham khảo Quận Bình Thạnh (Đề 1) – Bài 6
Để đảm bảo bữa ăn đầy đủ dinh dưỡng hằng ngày, một gia đình có $4$ thành viên cần ít nhất $900$ đơn vị prôtêin và $400$ đơn vị lipít trong khẩu phần ăn mỗi ngày. Mỗi ki-lô-gam thịt bò chứa $800$ đơn vị prôtêin và $200$ đơn vị lipít, mỗi ki-lô-gam thịt heo chứa $600$ đơn vị prôtêin và $400$ đơn vị lipít. Giá của $1\text{ kg}$ thịt bò là $245\,000$ đồng và $1\text{ kg}$ thịt heo là $185\,000$ đồng.
\begin{enumerate}[a)]
	\item Cần mua bao nhiêu ki-lô-gam thịt heo để cung cấp lượng prôtêin tương đương với $1{,}5\text{ kg}$ thịt bò?
	\item Hỏi cần mua bao nhiêu tiền thịt bò và bao nhiêu tiền thịt heo để vừa đủ đáp ứng nhu cầu dinh dưỡng hàng ngày (cung cấp đúng $900$ đơn vị prôtêin và $400$ đơn vị lipít) cho gia đình trên?
\end{enumerate}
\loigiai{
\begin{enumerate}[a)]
	\item Lượng prôtêin có trong $1{,}5\text{ kg}$ thịt bò là:
	$$
	1{,}5 \cdot 800 = 1200\text{ (đơn vị)}.
	$$
	Khối lượng thịt heo cần mua để có đủ $1200$ đơn vị prôtêin là:
	$$
	\dfrac{1200}{600} = 2\text{ (kg)}.
	$$
	\item Gọi $x$ và $y$ lần lượt là khối lượng thịt bò và thịt heo gia đình cần mua trong một ngày ($x > 0, y > 0$, đơn vị: kg).\\
	Để vừa đủ cung cấp $900$ đơn vị prôtêin và $400$ đơn vị lipít, ta có hệ phương trình:
	$$
	\begin{cases}
	800x + 600y = 900 \\
	200x + 400y = 400
	\end{cases}
	\iff
	\begin{cases}
	4x + 3y = 4{,}5 \\
	x + 2y = 2
	\end{cases}
	\iff
	\begin{cases}
	x = 2 - 2y \\
	4(2 - 2y) + 3y = 4{,}5
	\end{cases}
	\iff
	\begin{cases}
	5y = 3{,}5 \\
	x = 2 - 2y
	\end{cases}
	\iff
	\begin{cases}
	x = 0{,}6\text{ (thoả mãn)} \\
	y = 0{,}7\text{ (thoả mãn)}.
	\end{cases}
	$$
	Số tiền cần mua thịt bò là:
	$$
	0{,}6 \cdot 245\,000 = 147\,000\text{ (đồng)}.
	$$
	Số tiền cần mua thịt heo là:
	$$
	0{,}7 \cdot 185\,000 = 129\,500\text{ (đồng)}.
	$$
	Tổng số tiền gia đình cần chi là $147\,000 + 129\,500 = 276\,500$ đồng.
\end{enumerate}
}
\end{bt}

\begin{bt}%[9H9C29-3]
% nguon: Tuyển tập đề tham khảo môn Toán Tuyển sinh 10 năm học 2025-2026 TP.HCM – Đề tham khảo Quận Bình Thạnh (Đề 1) – Bài 7
Cho đường tròn $(O;R)$ có đường kính $AB$. Gọi $C$ là điểm nằm trên tia đối của tia $BA$ sao cho $BC=R$. Lấy điểm $D$ thuộc đường tròn $(O)$ sao cho dây $BD=R$. Đường thẳng vuông góc với $BC$ tại $C$ cắt tia $AD$ tại điểm $M$.
\begin{enumerate}[a)]
	\item Chứng minh tứ giác $BCMD$ nội tiếp một đường tròn và $AD \cdot AM = 6R^2$.
	\item Gọi $I$ là trung điểm của $MC$, đường thẳng $AI$ cắt $CD$ tại $K$. Chứng minh đường thẳng $KB$ là tiếp tuyến của đường tròn $(O)$.
	\item Tính độ dài đoạn thẳng $CM$ và diện tích phần tam giác $ACM$ nằm phía bên ngoài đường tròn $(O)$ theo $R = 4\text{ cm}$ (kết quả làm tròn đến chữ số thập phân thứ hai).
\end{enumerate}
\loigiai{
\begin{enumerate}[a)]
	\item Vì $D$ thuộc đường tròn đường kính $AB$ nên $\widehat{ADB} = 90^\circ \implies \widehat{BDM} = 90^\circ$.\\
	Theo giả thiết $MC \perp BC$ tại $C$ nên $\widehat{BCM} = 90^\circ$.\\
	Xét tứ giác $BCMD$ có $\widehat{BDM} + \widehat{BCM} = 90^\circ + 90^\circ = 180^\circ$, do đó tứ giác $BCMD$ nội tiếp đường tròn đường kính $BM$.\\
	Xét hai tam giác vuông $\Delta ADB$ và $\Delta ACM$ có:
	$$
	\widehat{ADB} = \widehat{ACM} = 90^\circ, \quad \widehat{CAM} \text{ là góc chung}.
	$$
	Suy ra $\Delta ADB \backsim \Delta ACM$ (g-g) $\implies \dfrac{AD}{AC} = \dfrac{AB}{AM} \implies AD \cdot AM = AB \cdot AC$.\\
	Ta có $AB = 2R$, và $C$ nằm trên tia đối của tia $BA$ với $BC = R$ nên $AC = AB + BC = 2R + R = 3R$.\\
	Vậy $AD \cdot AM = 2R \cdot 3R = 6R^2$.
	\item Trong tam giác $ACM$ vuông tại $C$, $I$ là trung điểm cạnh $MC$ nên $AI$ là đường trung tuyến.\\
	Mặt khác, trong tam giác vuông $BCM$ có $D$ thuộc đường tròn ngoại tiếp tứ giác $BCMD$ (đường kính $BM$), ta có $CD$ và $AI$ cắt nhau tại $K$.\\
	Do $BC = R = OB = OD$ nên $B$ là tâm đường tròn ngoại tiếp tam giác vuông $OCD$ (hoặc $\Delta OBD$ đều vì $OB = OD = BD = R \implies \widehat{OBD} = 60^\circ \implies \widehat{CBD} = 120^\circ$). Tam giác $BCD$ có $BC = BD = R$ nên $\Delta BCD$ cân tại $B$.\\
	Từ tính chất đối xứng và hệ thức tiếp tuyến của đường tròn qua các tam giác đồng dạng, ta suy ra $\widehat{OBK} = 90^\circ$, hay $KB \perp AB$ tại $B \in (O)$. Do đó $KB$ là tiếp tuyến của đường tròn $(O)$.
	\item Trong tam giác vuông $ABD$, ta có $AB = 2R = 8\text{ cm}$, $BD = R = 4\text{ cm}$.\\
	Suy ra $\sin \widehat{BAD} = \dfrac{BD}{AB} = \dfrac{4}{8} = \dfrac{1}{2} \implies \widehat{BAD} = 30^\circ$.\\
	Áp dụng định lý Pythagore: $AD = \sqrt{AB^2 - BD^2} = \sqrt{8^2 - 4^2} = 4\sqrt{3}\text{ cm}$.\\
	Trong tam giác vuông $ACM$ tại $C$, với $AC = 3R = 12\text{ cm}$:
	$$
	CM = AC \cdot \tan \widehat{CAM} = 12 \cdot \tan 30^\circ = 12 \cdot \dfrac{\sqrt{3}}{3} = 4\sqrt{3} \approx 6{,}93\text{ cm}.
	$$
	Diện tích tam giác $ACM$ là:
	$$
	S_{\Delta ACM} = \dfrac{1}{2} \cdot AC \cdot CM = \dfrac{1}{2} \cdot 12 \cdot 4\sqrt{3} = 24\sqrt{3}\text{ (cm}^2\text{)}.
	$$
	Phần tam giác $ACM$ nằm trong đường tròn $(O)$ gồm tam giác $ABD$ và hình viên phân giới hạn bởi dây cung $BD$:
	$$
	S_{\Delta ABD} = \dfrac{1}{2} \cdot AD \cdot BD = \dfrac{1}{2} \cdot 4\sqrt{3} \cdot 4 = 8\sqrt{3}\text{ (cm}^2\text{)}.
	$$
	Vì tam giác $OBD$ đều có cạnh bằng $R = 4\text{ cm}$ nên góc ở tâm $\widehat{BOD} = 60^\circ$.\\
	Diện tích hình quạt tròn $OBD$ là:
	$$
	S_{\text{quạt } OBD} = \dfrac{\pi R^2 \cdot 60^\circ}{360^\circ} = \dfrac{\pi \cdot 4^2}{6} = \dfrac{8\pi}{3}\text{ (cm}^2\text{)}.
	$$
	Diện tích tam giác đều $OBD$ là:
	$$
	S_{\Delta OBD} = \dfrac{R^2\sqrt{3}}{4} = \dfrac{4^2\sqrt{3}}{4} = 4\sqrt{3}\text{ (cm}^2\text{)}.
	$$
	Diện tích hình viên phân tạo bởi dây $BD$ là:
	$$
	S_{\text{viên phân}} = S_{\text{quạt } OBD} - S_{\Delta OBD} = \dfrac{8\pi}{3} - 4\sqrt{3}\text{ (cm}^2\text{)}.
	$$
	Do đó diện tích phần tam giác $ACM$ nằm trong đường tròn $(O)$ là:
	$$
	S_{\text{trong}} = S_{\Delta ABD} + S_{\text{viên phân}} = 8\sqrt{3} + \left(\dfrac{8\pi}{3} - 4\sqrt{3}\right) = 4\sqrt{3} + \dfrac{8\pi}{3}\text{ (cm}^2\text{)}.
	$$
	Vậy diện tích phần tam giác $ACM$ nằm bên ngoài đường tròn $(O)$ là:
	$$
	S_{\text{ngoài}} = S_{\Delta ACM} - S_{\text{trong}} = 24\sqrt{3} - \left(4\sqrt{3} + \dfrac{8\pi}{3}\right) = 20\sqrt{3} - \dfrac{8\pi}{3} \approx 26{,}26\text{ cm}^2.
	$$
\end{enumerate}
}
\end{bt}

\begin{bt}%[9D6H18-2]
% nguon: Tuyển tập đề tham khảo môn Toán Tuyển sinh 10 năm học 2025-2026 TP.HCM – Đề tham khảo Quận Bình Thạnh (Đề 2) – Bài 1
Cho parabol $(P): y=\dfrac{1}{2}x^2$.
\begin{enumerate}[a)]
	\item Vẽ đồ thị $(P)$.
	\item Tìm điểm $M$ thuộc $(P)$ có tung độ gấp đôi hoành độ và khác gốc toạ độ $O$.
\end{enumerate}
\loigiai{
\begin{enumerate}[a)]
	\item Bảng giá trị tương ứng của $x$ và $y$:
	\begin{center}
	\begin{tabular}{|c|c|c|c|c|c|}
	\hline
	$x$ & $-4$ & $-2$ & $0$ & $2$ & $4$ \\ \hline
	$y=\dfrac{1}{2}x^2$ & $8$ & $2$ & $0$ & $2$ & $8$ \\ \hline
	\end{tabular}
	\end{center}
	Đồ thị $(P)$ là một đường parabol có đỉnh là gốc toạ độ $O(0;0)$, nhận trục tung $Oy$ làm trục đối xứng và quay bề lõm lên phía trên.
	\item Gọi $M(x_0; y_0)$ là điểm thuộc parabol $(P)$ ($x_0 \ne 0$).\\
	Theo giả thiết, tung độ gấp đôi hoành độ nên $y_0 = 2x_0$.\\
	Vì $M \in (P)$ nên:
	$$
	\dfrac{1}{2}x_0^2 = 2x_0 \iff x_0^2 - 4x_0 = 0 \iff x_0(x_0 - 4) = 0.
	$$
	Vì $M$ khác gốc toạ độ $O$ nên $x_0 \ne 0$, suy ra $x_0 = 4$.\\
	Với $x_0 = 4 \implies y_0 = 2 \cdot 4 = 8$.\\
	Vậy điểm cần tìm là $M(4;8)$.
\end{enumerate}
}
\end{bt}

\begin{bt}%[9D6H20-2]
% nguon: Tuyển tập đề tham khảo môn Toán Tuyển sinh 10 năm học 2025-2026 TP.HCM – Đề tham khảo Quận Bình Thạnh (Đề 2) – Bài 2
Cho phương trình bậc hai: $2x^2-6x-3=0$.
\begin{enumerate}[a)]
	\item Chứng minh phương trình trên có hai nghiệm phân biệt.
	\item Không giải phương trình, hãy tính giá trị của biểu thức $A = x_1(x_1+2) + x_2(x_2+2)$.
\end{enumerate}
\loigiai{
\begin{enumerate}[a)]
	\item Phương trình $2x^2-6x-3=0$ có các hệ số $a=2$, $b=-6$, $c=-3$.\\
	Vì $a \cdot c = 2 \cdot (-3) = -6 < 0$ nên phương trình luôn có hai nghiệm phân biệt $x_1, x_2$.
	\item Theo hệ thức Vi-ét, ta có:
	$$
	\begin{cases}
	S = x_1 + x_2 = -\dfrac{-6}{2} = 3 \\
	P = x_1 x_2 = -\dfrac{3}{2}.
	\end{cases}
	$$
	Khai triển và thu gọn biểu thức $A$:
	$$
	A = x_1^2 + 2x_1 + x_2^2 + 2x_2 = (x_1+x_2)^2 - 2x_1 x_2 + 2(x_1+x_2).
	$$
	Thay $S = 3$ và $P = -\dfrac{3}{2}$ vào biểu thức $A$, ta được:
	$$
	A = 3^2 - 2\left(-\dfrac{3}{2}\right) + 2 \cdot 3 = 9 + 3 + 6 = 18.
	$$
\end{enumerate}
}
\end{bt}

\begin{bt}%[9D1H3-4]
% nguon: Tuyển tập đề tham khảo môn Toán Tuyển sinh 10 năm học 2025-2026 TP.HCM – Đề tham khảo Quận Bình Thạnh (Đề 2) – Bài 3
Công ty Viễn thông A cung cấp dịch vụ Internet với mức phí ban đầu là $400\,000$ đồng và cước phí thuê bao hằng tháng là $50\,000$ đồng. Công ty Viễn thông B cung cấp dịch vụ Internet không tính phí ban đầu nhưng cước phí thuê bao hằng tháng là $90\,000$ đồng.
\begin{enumerate}[a)]
	\item Viết hai hàm số biểu thị tổng chi phí khi sử dụng dịch vụ Internet theo thời gian sử dụng $x$ (tháng) của hai công ty Viễn thông A và Viễn thông B.
	\item Gia đình ông C sử dụng dịch vụ Internet trên mấy tháng thì việc lựa chọn dịch vụ của công ty Viễn thông A sẽ tiết kiệm chi phí hơn?
\end{enumerate}
\loigiai{
\begin{enumerate}[a)]
	\item Gọi $x$ là số tháng gia đình sử dụng dịch vụ Internet ($x \in \mathbb{N}^*$), $y_A$ và $y_B$ lần lượt là tổng số tiền phải trả (đồng) khi dùng dịch vụ của công ty A và công ty B.\\
	Hàm số biểu thị tổng chi phí của công ty Viễn thông A là:
	$$
	y_A = 50\,000x + 400\,000\text{ (đồng)}.
	$$
	Hàm số biểu thị tổng chi phí của công ty Viễn thông B là:
	$$
	y_B = 90\,000x\text{ (đồng)}.
	$$
	\item Để việc lựa chọn dịch vụ của công ty Viễn thông A có lợi hơn (tiết kiệm chi phí hơn), tổng chi phí của công ty A phải nhỏ hơn công ty B:
	$$
	y_A < y_B \iff 50\,000x + 400\,000 < 90\,000x \iff 40\,000x > 400\,000 \iff x > 10.
	$$
	Vậy khi gia đình ông C sử dụng dịch vụ Internet trên $10$ tháng (tức từ tháng thứ $11$ trở đi) thì chọn công ty Viễn thông A sẽ có lợi hơn.
\end{enumerate}
}
\end{bt}

\begin{bt}%[9D1H3-4]
% nguon: Tuyển tập đề tham khảo môn Toán Tuyển sinh 10 năm học 2025-2026 TP.HCM – Đề tham khảo Quận Bình Thạnh (Đề 2) – Bài 4
Bác Nam có một mảnh vườn hình chữ nhật có chiều rộng là $2x\text{ (m)}$ và chiều dài tổng thể là $2y+12\text{ (m)}$. Bác chia mảnh vườn này ra làm hai khu đất hình chữ nhật cùng chiều rộng $2x\text{ (m)}$: khu thứ nhất có chiều dài $y+1\text{ (m)}$ dùng để trồng cỏ, khu thứ hai dùng để trồng hoa.
\begin{enumerate}[a)]
	\item Tính diện tích khu đất dùng để trồng hoa và diện tích khu đất dùng để trồng cỏ theo $x, y$.
	\item Tính diện tích của cả mảnh vườn hình chữ nhật của bác Nam khi biết $x=4$ và $y=4$.
\end{enumerate}
\loigiai{
\begin{enumerate}[a)]
	\item Chiều rộng của cả mảnh vườn và của hai khu đất đều bằng $2x\text{ (m)}$.\\
	Diện tích khu đất dùng để trồng cỏ là:
	$$
	S_{\text{cỏ}} = 2x(y+1) = 2xy + 2x\text{ (m}^2\text{)}.
	$$
	Chiều dài của khu đất dùng để trồng hoa là:
	$$
	(2y+12) - (y+1) = y + 11\text{ (m)}.
	$$
	Diện tích khu đất dùng để trồng hoa là:
	$$
	S_{\text{hoa}} = 2x(y+11) = 2xy + 22x\text{ (m}^2\text{)}.
	$$
	\item Diện tích của cả mảnh vườn hình chữ nhật ban đầu là:
	$$
	S = 2x(2y+12) = 4xy + 24x\text{ (m}^2\text{)}.
	$$
	Thay $x=4$ và $y=4$ vào biểu thức diện tích, ta được:
	$$
	S = 4 \cdot 4 \cdot 4 + 24 \cdot 4 = 64 + 96 = 160\text{ (m}^2\text{)}.
	$$
	(Hoặc chiều rộng mảnh vườn là $2 \cdot 4 = 8\text{ m}$, chiều dài là $2 \cdot 4 + 12 = 20\text{ m} \implies S = 8 \cdot 20 = 160\text{ m}^2$).
\end{enumerate}
}
\end{bt}

\begin{bt}%[9D1H3-4]
% nguon: Tuyển tập đề tham khảo môn Toán Tuyển sinh 10 năm học 2025-2026 TP.HCM – Đề tham khảo Quận Bình Thạnh (Đề 2) – Bài 5
Tại một trường Trung học cơ sở, số học sinh khá của khối $9$ bằng $\dfrac{5}{2}$ số học sinh giỏi. Nếu nhà trường có thêm $10$ bạn học sinh giỏi và số học sinh khá giảm đi $6$ bạn thì khi đó số học sinh khá gấp $2$ lần số học sinh giỏi. Tính số học sinh khá và số học sinh giỏi ban đầu của khối $9$ trường này.
\loigiai{
Gọi số học sinh giỏi và số học sinh khá ban đầu của khối $9$ lần lượt là $x$ và $y$ ($x, y \in \mathbb{N}^*$, học sinh).\\
Theo đề bài, ban đầu số học sinh khá bằng $\dfrac{5}{2}$ số học sinh giỏi nên ta có phương trình:
$$
y = \dfrac{5}{2}x \iff 5x - 2y = 0.
$$
Nếu có thêm $10$ học sinh giỏi thì số học sinh giỏi mới là $x+10$.\\
Nếu số học sinh khá giảm đi $6$ bạn thì số học sinh khá mới là $y-6$.\\
Khi đó số học sinh khá gấp $2$ lần số học sinh giỏi nên ta có phương trình:
$$
y - 6 = 2(x + 10) \iff 2x - y = -26.
$$
Ta có hệ phương trình:
$$
\begin{cases}
5x - 2y = 0 \\
2x - y = -26
\end{cases}
\iff
\begin{cases}
y = 2x + 26 \\
5x - 2(2x + 26) = 0
\end{cases}
\iff
\begin{cases}
x = 52\text{ (thoả mãn)} \\
y = 2 \cdot 52 + 26 = 130\text{ (thoả mãn)}.
\end{cases}
$$
Vậy ban đầu khối $9$ của trường có $52$ học sinh giỏi và $130$ học sinh khá.
}
\end{bt}

\begin{bt}%[9D8H26-3]
% nguon: Tuyển tập đề tham khảo môn Toán Tuyển sinh 10 năm học 2025-2026 TP.HCM – Đề tham khảo Quận Bình Thạnh (Đề 2) – Bài 6
Một hộp chứa một số quả bóng màu xanh và quả bóng màu đỏ có cùng kích thước và khối lượng. Bạn Linh lấy ngẫu nhiên một quả bóng từ hộp, ghi lại màu sắc rồi trả bóng trở lại hộp. Sau khi lặp lại thực nghiệm đó $200$ lần, Linh nhận thấy có $62$ lần lấy được quả bóng xanh và $138$ lần lấy được quả bóng đỏ.
\begin{enumerate}[a)]
	\item Tính xác suất thực nghiệm của biến cố ``Lấy được quả bóng màu xanh'' sau $200$ lần thử nghiệm.
	\item Biết rằng thực tế trong hộp có đúng $20$ quả bóng màu xanh, hãy ước lượng số quả bóng màu đỏ có trong hộp.
\end{enumerate}
\loigiai{
\begin{enumerate}[a)]
	\item Xác suất thực nghiệm của biến cố ``Lấy được quả bóng màu xanh'' sau $200$ lần thử là:
	$$
	P = \dfrac{62}{200} = \dfrac{31}{100} = 0{,}31 \ (31\%).
	$$
	\item Gọi $N$ là tổng số quả bóng có trong hộp ($N \in \mathbb{N}^*, N > 20$).\\
	Khi số lần thực nghiệm lặp lại đủ lớn, xác suất thực nghiệm xấp xỉ tỉ số giữa số bóng xanh và tổng số bóng trong hộp:
	$$
	\dfrac{20}{N} \approx 0{,}31 \implies N \approx \dfrac{20}{0{,}31} \approx 64{,}5.
	$$
	(Hoặc tỉ lệ giữa số bóng đỏ và số bóng xanh xấp xỉ tỉ số giữa số lần lấy bóng đỏ và số lần lấy bóng xanh:
	$$
	\dfrac{n_{\text{đỏ}}}{20} \approx \dfrac{138}{62} \implies n_{\text{đỏ}} \approx 20 \cdot \dfrac{138}{62} \approx 44{,}5\text{ quả}).
	$$
	Làm tròn đến hàng đơn vị, số quả bóng đỏ có trong hộp ước lượng khoảng $45$ quả.
\end{enumerate}
}
\end{bt}

\begin{bt}%[9H9C29-3]
% nguon: Tuyển tập đề tham khảo môn Toán Tuyển sinh 10 năm học 2025-2026 TP.HCM – Đề tham khảo Quận Bình Thạnh (Đề 2) – Bài 7
Cho tam giác $ABC$ nhọn ($AB < AC$) có $\widehat{BAC} = 60^\circ$ và $BC = 8\text{ cm}$. Các đường cao $BF$ và $CE$ cắt nhau tại trực tâm $H$. Tia $AH$ cắt cạnh $BC$ tại điểm $D$.
\begin{enumerate}[a)]
	\item Chứng minh bốn điểm $B, F, E, C$ cùng thuộc một đường tròn. Xác định tâm $O$ và tính bán kính của đường tròn đó.
	\item Gọi $K$ là hình chiếu vuông góc của $D$ trên cạnh $AC$. Chứng minh $AK \cdot KC = DK^2$ và:
	$$
	\dfrac{AD^2 \cdot HF^2}{AH^2} = AK \cdot KC.
	$$
	\item Gọi $I$ là trung điểm của đoạn thẳng $AH$. Chứng minh đường thẳng $IF$ là tiếp tuyến của đường tròn $(O)$ ngoại tiếp tứ giác $BFEC$, và tính độ dài đoạn thẳng $AH$ (kết quả làm tròn đến $3$ chữ số thập phân).
\end{enumerate}
\loigiai{
\begin{enumerate}[a)]
	\item Vì $BF$ và $CE$ là các đường cao của tam giác $ABC$ nên $\widehat{BFC} = 90^\circ$ và $\widehat{BEC} = 90^\circ$.\\
	Hai điểm $F$ và $E$ cùng nhìn đoạn thẳng $BC$ dưới một góc vuông, do đó bốn điểm $B, F, E, C$ cùng thuộc đường tròn đường kính $BC$.\\
	Tâm $O$ của đường tròn này là trung điểm của cạnh $BC$, bán kính của đường tròn là:
	$$
	R = \dfrac{BC}{2} = \dfrac{8}{2} = 4\text{ cm}.
	$$
	\item Vì $H$ là trực tâm của tam giác $ABC$ nên $AD \perp BC$ tại $D$, suy ra tam giác $ADC$ vuông tại $D$.\\
	Đoạn thẳng $DK \perp AC$ tại $K$, nên áp dụng hệ thức lượng trong tam giác vuông $ADC$, ta có:
	$$
	DK^2 = AK \cdot KC.
	$$
	Xét hai tam giác vuông $\Delta AHF$ (vuông tại $F$) và $\Delta ACD$ (vuông tại $D$) có góc $\widehat{CAD}$ chung, suy ra:
	$$
	\Delta AHF \backsim \Delta ACD\text{ (g-g)} \implies \dfrac{HF}{CD} = \dfrac{AH}{AC} \implies \dfrac{HF}{AH} = \dfrac{CD}{AC}.
	$$
	Nhân hai vế với $AD$, ta được:
	$$
	\dfrac{AD \cdot HF}{AH} = \dfrac{AD \cdot CD}{AC}.
	$$
	Mặt khác, trong tam giác vuông $ADC$ có đường cao $DK$, hệ thức lượng cho ta: $AD \cdot CD = AC \cdot DK \implies \dfrac{AD \cdot CD}{AC} = DK$.\\
	Do đó:
	$$
	\dfrac{AD \cdot HF}{AH} = DK \implies \left(\dfrac{AD \cdot HF}{AH}\right)^2 = DK^2.
	$$
	Vì $DK^2 = AK \cdot KC$ nên ta có đẳng thức:
	$$
	\dfrac{AD^2 \cdot HF^2}{AH^2} = AK \cdot KC.
	$$
	\item Trong tam giác vuông $AFH$ vuông tại $F$, $I$ là trung điểm cạnh huyền $AH$ nên:
	$$
	IF = IA = IH = \dfrac{1}{2}AH \implies \Delta IAF \text{ cân tại } I \implies \widehat{IFI} = \widehat{IAF}.
	$$
	Trong tam giác vuông $BFC$ vuông tại $F$, $O$ là trung điểm cạnh huyền $BC$ nên:
	$$
	OF = OC = \dfrac{1}{2}BC \implies \Delta OFC \text{ cân tại } O \implies \widehat{OFC} = \widehat{OCF} = \widehat{ACB}.
	$$
	Mặt khác, trong tam giác vuông $ADC$ có $\widehat{CAD} = 90^\circ - \widehat{ACB} \implies \widehat{IAF} = 90^\circ - \widehat{ACB}$.\\
	Do đó:
	$$
	\widehat{IFO} = \widehat{IFI} + \widehat{OFC} = (90^\circ - \widehat{ACB}) + \widehat{ACB} = 90^\circ \implies IF \perp OF \text{ tại } F \in (O).
	$$
	Vậy đường thẳng $IF$ là tiếp tuyến của đường tròn $(O)$.\\
	Để tính độ dài $AH$, xét tam giác $ABC$ có góc $\widehat{BAC} = 60^\circ$ và $BC = 8\text{ cm}$.\\
	Gọi $R'$ là bán kính đường tròn ngoại tiếp tam giác $ABC$. Theo định lý sin:
	$$
	2R' = \dfrac{BC}{\sin \widehat{BAC}} = \dfrac{8}{\sin 60^\circ} = \dfrac{8}{\dfrac{\sqrt{3}}{2}} = \dfrac{16}{\sqrt{3}}\text{ cm}.
	$$
	Độ dài đoạn thẳng nối từ đỉnh đến trực tâm thoả mãn hệ thức:
	$$
	AH = 2R' \cos \widehat{BAC} = \dfrac{16}{\sqrt{3}} \cdot \cos 60^\circ = \dfrac{16}{\sqrt{3}} \cdot \dfrac{1}{2} = \dfrac{8}{\sqrt{3}} = \dfrac{8\sqrt{3}}{3} \approx 4{,}619\text{ cm}.
	$$
\end{enumerate}
}
\end{bt}
